\documentclass[12pt,a4paper]{article}

% ── Packages ──
\usepackage[utf8]{inputenc}
\usepackage[T1]{fontenc}
\usepackage{amsmath,amssymb,amsfonts}
\usepackage{graphicx}
\usepackage{booktabs}
\usepackage{hyperref}
\usepackage[margin=2.5cm]{geometry}
\usepackage{natbib}
\usepackage{caption}
\usepackage{subcaption}
\usepackage{multirow}
\usepackage{array}
\usepackage{xcolor}
\usepackage{algorithm}
\usepackage{algorithmic}

\hypersetup{colorlinks=true, linkcolor=blue!60!black, citecolor=blue!60!black, urlcolor=blue!60!black}

\newcolumntype{R}{>{\raggedleft\arraybackslash}p{1.1cm}}

\title{When Should Actuaries Trust Neural Networks?\\
\large A Comprehensive Benchmark of Classical and Neural Mortality Models\\with a Hybrid LC-ResNet and Practical Decision Framework}

\author{Amine Manai\\
\textit{Master 1 Actuariat, Le Mans Universit\'e}\\
\texttt{amine.manai@esprit.tn}\\[6pt]
\small July 2026
}

\date{}

\begin{document}

\maketitle

% ══════════════════════════════════════════════════════════════
\begin{abstract}
We present a comprehensive benchmark comparing 15 mortality forecasting models---8~classical (implemented from scratch), 6~neural network architectures, and an original hybrid---across 8~countries from the Human Mortality Database (1950--2023). Using rolling-origin validation with actuarial metrics (life expectancy $e_0$, annuity values $\ddot{a}_{65}$), we address practical questions that existing studies leave unanswered: which model to use with short historical data, which is most robust to mortality shocks (COVID-19), and when the added complexity of neural networks is justified.

Our main contributions are: (1)~a practical decision framework mapping data characteristics to optimal model choice; (2)~\textbf{LC-ResNet}, a hybrid model combining a Poisson Lee-Carter skeleton with a neural residual correction subject to horizon-dependent shrinkage $\exp(-\lambda h)$, ensuring stable long-horizon extrapolation while capturing short-term non-linearities; and (3)~an actuarial case study quantifying model risk in EUR on a French annuity portfolio under Solvency~II.

Key findings: the random walk with drift remains a strong baseline at all horizons; neural $\kappa_t$ replacements (LSTM, GRU) outperform classical Lee-Carter at short horizons ($h \le 5$) but offer diminishing gains beyond; LC-ResNet ranks among the top models for young ages while preserving interpretability; and the inter-model spread on a 1,000-life annuity portfolio reaches 6.3~million~EUR, underscoring the materiality of model risk.

\medskip
\noindent\textbf{Keywords:} mortality forecasting, Lee-Carter, neural networks, LSTM, model risk, longevity risk, Solvency~II, annuity pricing, decision framework.

\medskip
\noindent\textbf{JEL:} C45, C52, C53, G22, J11.
\end{abstract}

\newpage
\tableofcontents
\newpage

% ══════════════════════════════════════════════════════════════
\section{Introduction}
\label{sec:intro}

Longevity risk---the risk that policyholders live longer than expected---is a primary concern for life insurers and pension funds. Under Solvency~II, the capital requirement for longevity risk (SCR\textsubscript{longevity}) typically represents 5--15\% of the best estimate of annuity provisions, making the choice of mortality model a capital allocation decision rather than a purely academic exercise.

The Lee-Carter family of models \citep{lee1992modeling} has dominated actuarial practice for over three decades. In parallel, neural network approaches have shown promising results in recent literature \citep{richman2019neural, nigri2019deep, perla2021brief}. Yet no study provides a practical guide for actuaries to choose between classical and neural models based on their specific context: data availability, forecast horizon, age segment of interest, and regulatory constraints.

This paper fills that gap by benchmarking 15~models across 8~countries with rolling-origin validation, then distilling the results into a decision framework and quantifying the financial consequences through an actuarial case study.

\subsection{Contributions}

\begin{enumerate}
    \item A \textbf{comprehensive benchmark} of 15~models (8~classical, 6~neural, 1~hybrid) on 8~countries from the Human Mortality Database, evaluated with both statistical and actuarial metrics under rolling-origin validation.
    \item \textbf{LC-ResNet}, an original hybrid model combining a Poisson Lee-Carter skeleton with a neural residual network and horizon-dependent shrinkage, designed to preserve interpretability while capturing short-term non-linearities.
    \item A \textbf{practical decision framework} mapping data characteristics (history length, forecast horizon, age segment, interpretability needs) to optimal model choice, based on scenario-specific evaluations.
    \item An \textbf{actuarial case study} pricing a French annuity portfolio and quantifying model risk in EUR, including a Solvency~II longevity shock analysis.
\end{enumerate}

\subsection{Related work}

\citet{lee1992modeling} introduced the seminal Lee-Carter model; subsequent extensions include Lee-Miller \citep{lee2001evaluating}, Booth-Maindonald-Smith \citep{booth2002applying}, the Poisson formulation of \citet{brouhns2002poisson}, the two-factor CBD model for older ages \citep{cairns2006two}, and the functional data approach of \citet{hyndman2007robust}. On the neural side, \citet{nigri2019deep} applied LSTMs to the $\kappa_t$ index, while \citet{richman2019neural} proposed a feed-forward network with entity embeddings for multi-population mortality. \citet{perla2021brief} provided a systematic review of deep learning methods for mortality forecasting.

Our work differs by (i)~benchmarking all these approaches under identical conditions with rolling-origin validation, (ii)~evaluating on actuarial rather than purely statistical metrics, and (iii)~proposing a hybrid model and a practical decision framework.

% ══════════════════════════════════════════════════════════════
\section{Data}
\label{sec:data}

We use data from the \textbf{Human Mortality Database} (HMD; \url{www.mortality.org}) for 8~countries: France (FRATNP), England \& Wales (GBRTENW), United States (USA), Japan (JPN), Italy (ITA), Spain (ESP), Sweden (SWE), and Netherlands (NLD). For each country, we extract $1 \times 1$ (single-year age, single calendar year) central death rates $m_{x,t}$ and exposures $E_{x,t}$ for ages 0--100 and years 1950--2023, total population (both sexes combined).

The period 1950--2023 is chosen to include sufficient historical depth for long-horizon evaluation and to encompass the COVID-19 mortality shock (2020--2023), which serves as a natural out-of-sample stress test.

Zero or missing rates are floored at $10^{-6}$ before taking logarithms.

% ══════════════════════════════════════════════════════════════
\section{Models}
\label{sec:models}

All classical models are implemented from scratch in Python using NumPy and SciPy; neural models use PyTorch. All share a common \texttt{MortalityModel} interface (\texttt{fit}, \texttt{forecast}, \texttt{simulate}) to ensure evaluation consistency.

\subsection{Classical models}

\paragraph{Lee-Carter (LC).} \citet{lee1992modeling}: $\ln m_{x,t} = a_x + b_x \kappa_t + \varepsilon_{x,t}$, estimated by SVD with identifiability constraints $\sum_x b_x = 1$, $\sum_t \kappa_t = 0$, and $\kappa_t$ re-estimated on observed deaths via constrained optimization. Forecast: $\kappa_t$ follows a random walk with drift (RWD).

\paragraph{Lee-Miller (LM).} \citet{lee2001evaluating}: extends LC by adjusting $\kappa_t$ so that the model-implied $e_0$ matches the observed $e_0$ each year, solved via Brent's root-finding algorithm.

\paragraph{Booth-Maindonald-Smith (BMS).} \citet{booth2002applying}: selects the optimal fitting period by testing the linearity of $\kappa_t$ over sub-periods, then applies standard LC on the selected window.

\paragraph{Poisson Lee-Carter (PLC).} \citet{brouhns2002poisson}: $D_{x,t} \sim \text{Poisson}(E_{x,t} \cdot \exp(a_x + b_x \kappa_t))$, estimated by maximum likelihood via alternating unidimensional Newton-Raphson iterations.

\paragraph{Cairns-Blake-Dowd (CBD).} \citet{cairns2006two}: a two-factor model for ages 60--100: $\text{logit}\, q_{x,t} = \kappa_t^{(1)} + \kappa_t^{(2)}(x - \bar{x})$, with bivariate RWD on $(\kappa_t^{(1)}, \kappa_t^{(2)})$.

\paragraph{Hyndman-Ullah (HU).} \citet{hyndman2007robust}: functional data analysis with penalized B-spline smoothing, $K$-component functional PCA, and ARIMA forecasting on each score series. We use $K=6$ components.

\paragraph{Baselines.} Random walk with drift (RWD) on age-specific $\ln m_{x,t}$, and frozen rates (last observed year projected forward).

\subsection{Neural models}

All neural models share identical training hyperparameters for benchmark fairness: 200~epochs, patience~20 (early stopping on validation loss), learning rate $10^{-3}$ (Adam, weight decay $10^{-4}$), seed~42.

\paragraph{Neural $\kappa_t$ models (LSTM, GRU, Bi-LSTM, Transformer).} The Lee-Carter skeleton ($a_x$, $b_x$) is preserved; only the $\kappa_t$ forecast is replaced by a recurrent or attention-based architecture (2~layers, hidden size~64, lookback~20~years). This retains the interpretable age structure while potentially capturing non-linear temporal patterns.

\paragraph{FFNN with embeddings.} Inspired by \citet{richman2019neural}: an age embedding (dimension~8) concatenated with a continuous normalized year feature, fed through a 3-layer MLP (128--64--32). Outputs $\ln m_{x,t}$ directly.

\paragraph{CNN on the mortality surface.} Treats the $(\text{age} \times \text{year})$ matrix as a 2D image; three convolutional layers (16, 32, 64~filters, kernel size~3) with adaptive average pooling predict the next year's mortality profile from a sliding window.

\subsection{LC-ResNet: a hybrid model}
\label{sec:lcresnet}

Our original contribution combines a Poisson Lee-Carter skeleton with a neural residual correction:

\begin{enumerate}
    \item \textbf{Fit} a Poisson Lee-Carter to obtain $\hat{a}_x$, $\hat{b}_x$, $\hat{\kappa}_t$.
    \item \textbf{Compute structured residuals}: $r_{x,t} = \ln m_{x,t} - (\hat{a}_x + \hat{b}_x \hat{\kappa}_t)$.
    \item \textbf{Train a residual network} $f_\theta$ on features $[\tilde{x}, \tilde{t}, \tilde{b}_x, \tilde{\kappa}_t]$ (normalized) with architecture MLP(64, 32) to predict $r_{x,t}$.
    \item \textbf{Forecast with shrinkage}: for horizon $h$,
    \begin{equation}
        \widehat{\ln m}_{x,t+h} = \underbrace{\hat{a}_x + \hat{b}_x \hat{\kappa}_{t+h}}_{\text{LC skeleton}} + \underbrace{e^{-\lambda h} \cdot f_\theta(\tilde{x}, \widetilde{t+h}, \tilde{b}_x, \widetilde{\kappa}_{t+h})}_{\text{shrunk neural correction}}
        \label{eq:lcresnet}
    \end{equation}
\end{enumerate}

The shrinkage factor $e^{-\lambda h}$ (with $\lambda = 0.1$) is the key design choice: at $h=1$, $e^{-0.1} \approx 0.90$, so the neural correction is nearly fully active; at $h=20$, $e^{-2} \approx 0.14$, so the model reverts to stable LC extrapolation. This prevents the common failure mode of neural mortality models---extrapolation instability at long horizons---while preserving short-term flexibility.

The parameters $a_x$, $b_x$, $\kappa_t$ remain fully interpretable and accessible for actuarial reporting, unlike purely neural approaches.

% ══════════════════════════════════════════════════════════════
\section{Evaluation Protocol}
\label{sec:evaluation}

\subsection{Rolling-origin validation}

We use rolling-origin cross-validation adapted for time series: the model is trained on years $[1950, T_0]$, then evaluated on the subsequent $h$~years. The origin $T_0$ ranges from 1990 to 2018 in steps of~2, and horizons $h \in \{1, 5, 10, 20\}$~years. The model is fully re-trained at each origin---no incremental learning.

\subsection{Metrics}

We report both statistical and actuarial metrics:

\begin{itemize}
    \item \textbf{RMSE on $\ln m_{x,t}$}: root mean squared error on log central death rates.
    \item \textbf{RMSE on $e_0$}: error on period life expectancy at birth (in years).
    \item \textbf{RMSE on $\ddot{a}_{65}$}: error on the whole-life annuity-due factor at age~65 (interest rate $i = 2\%$).
\end{itemize}

Actuarial metrics are computed only when the model covers the full age grid (0--100); the CBD model, restricted to ages 60--100, is evaluated on $\ln m_{x,t}$ only.

\subsection{Statistical testing}

We use the Diebold-Mariano test \citep{diebold1995comparing} with the Harvey-Leybourne-Newbold correction \citep{harvey1997testing} for multi-step autocorrelation and small-sample adjustment, referencing a Student's $t$ distribution.

\subsection{Scenario-specific evaluations}

\begin{itemize}
    \item \textbf{Short history}: training truncated to 20, 30, or 50~years to assess data-efficiency.
    \item \textbf{Mortality shock}: training ends in 2019, evaluation on 2020--2023 (COVID-19).
    \item \textbf{Age groups}: separate evaluation on young (0--19), working (20--64), and elderly (65--100) ages.
\end{itemize}

% ══════════════════════════════════════════════════════════════
\section{Results}
\label{sec:results}

The full benchmark produces 29{,}536 observations across 15~models, 8~countries, 15~rolling origins, 7~horizons, and multiple metrics. We report cross-country averages; country-level results are available in the online repository.

\subsection{Main results: RMSE on $\ln m_{x,t}$}

Table~\ref{tab:rmse_logmx} reports the average RMSE on $\ln m_{x,t}$ by model and horizon.

\begin{table}[htbp]
\centering
\caption{RMSE on $\ln m_{x,t}$, averaged across 8 countries (lower is better). Best in \textbf{bold}.}
\label{tab:rmse_logmx}
\small
\begin{tabular}{l RRRR}
\toprule
Model & $h=1$ & $h=5$ & $h=10$ & $h=20$ \\
\midrule
\multicolumn{5}{l}{\textit{Classical}} \\
Lee-Carter       & 0.157 & 0.207 & 0.199 & 0.248 \\
Lee-Miller       & 0.146 & 0.196 & 0.187 & \textbf{0.232} \\
BMS              & 0.157 & 0.207 & 0.199 & 0.248 \\
Poisson LC       & 0.193 & 0.244 & 0.231 & 0.272 \\
CBD (60+ only)   & \textbf{0.081} & \textbf{0.126} & \textbf{0.095} & \textbf{0.120} \\
Hyndman-Ullah    & 0.318 & 0.351 & 0.329 & 0.355 \\
\midrule
\multicolumn{5}{l}{\textit{Baselines}} \\
Random walk      & 0.103 & 0.155 & 0.146 & 0.203 \\
Frozen rates     & 0.103 & 0.164 & 0.186 & 0.314 \\
\midrule
\multicolumn{5}{l}{\textit{Neural}} \\
LSTM $\kappa_t$  & 0.132 & 0.182 & 0.180 & 0.263 \\
GRU $\kappa_t$   & 0.132 & 0.183 & 0.182 & 0.274 \\
Bi-LSTM $\kappa_t$ & 0.133 & 0.188 & 0.195 & 0.288 \\
Transformer $\kappa_t$ & 0.141 & 0.207 & 0.225 & 0.327 \\
FFNN embeddings  & 0.233 & 0.286 & 0.290 & 0.368 \\
CNN surface      & 0.151 & 0.221 & 0.240 & 0.448 \\
\midrule
\multicolumn{5}{l}{\textit{Hybrid}} \\
\textbf{LC-ResNet} & 0.122 & 0.182 & 0.184 & 0.247 \\
\bottomrule
\end{tabular}
\\[4pt]
\footnotesize Note: CBD is evaluated only on ages 60--100, making its RMSE not directly comparable to full-grid models.
\end{table}

\paragraph{Key observations.}

\begin{enumerate}
    \item The \textbf{random walk with drift is a remarkably strong baseline}, outperforming most classical LC variants and several neural models at all horizons.
    \item \textbf{Neural $\kappa_t$ models (LSTM, GRU)} improve over classical Lee-Carter at short horizons ($h \le 5$) but offer diminishing or negative gains at $h = 20$.
    \item \textbf{LC-ResNet} achieves the best non-CBD RMSE at $h = 1$ (0.122) and remains competitive at all horizons. At $h = 20$, it matches Lee-Miller (0.247 vs.\ 0.232), demonstrating the effectiveness of the shrinkage mechanism.
    \item The \textbf{Transformer and CNN} are the weakest neural architectures, suffering from extrapolation instability at long horizons (RMSE 0.327 and 0.448 at $h=20$).
    \item \textbf{Hyndman-Ullah} and \textbf{FFNN} underperform across all horizons in this benchmark.
\end{enumerate}

\subsection{Actuarial metrics}

Table~\ref{tab:actuarial} reports RMSE on $e_0$ (in years) and on $\ddot{a}_{65}$ (dimensionless).

\begin{table}[htbp]
\centering
\caption{Actuarial metrics, averaged across 8 countries.}
\label{tab:actuarial}
\small
\begin{tabular}{l RR RR}
\toprule
& \multicolumn{2}{c}{RMSE $e_0$ (years)} & \multicolumn{2}{c}{RMSE $\ddot{a}_{65}$} \\
\cmidrule(lr){2-3} \cmidrule(lr){4-5}
Model & $h=5$ & $h=20$ & $h=5$ & $h=20$ \\
\midrule
Lee-Carter       & 0.495 & 0.832 & 0.212 & 0.397 \\
Lee-Miller       & \textbf{0.386} & \textbf{0.742} & 0.228 & 0.398 \\
Poisson LC       & 0.444 & 0.826 & 0.215 & 0.403 \\
Hyndman-Ullah    & 0.486 & 0.694 & 0.221 & 0.437 \\
Random walk      & 0.383 & 0.764 & \textbf{0.192} & \textbf{0.395} \\
LSTM $\kappa_t$  & 0.807 & 2.059 & 0.432 & 1.017 \\
GRU $\kappa_t$   & 0.815 & 2.181 & 0.429 & 1.065 \\
\textbf{LC-ResNet} & 0.421 & 0.846 & 0.207 & 0.403 \\
\bottomrule
\end{tabular}
\end{table}

A striking finding is that \textbf{neural $\kappa_t$ models degrade sharply on actuarial metrics} at long horizons: the LSTM's $e_0$ RMSE at $h=20$ is 2.06~years vs.\ 0.74 for Lee-Miller. This is because small systematic biases in $\kappa_t$ forecasts compound through the life table. LC-ResNet avoids this trap thanks to its shrinkage mechanism: at $h=20$, it performs comparably to the classical LC family on $e_0$ and $\ddot{a}_{65}$.

\subsection{Scenario results}

\paragraph{Short history.} Table~\ref{tab:short_history} shows that with only 20~years of training data, model differences narrow considerably. Frozen rates and LC-ResNet perform best; neural $\kappa_t$ models converge to the LC baseline because their recurrent networks lack sufficient training sequences.

\begin{table}[htbp]
\centering
\caption{RMSE $\ln m_{x,t}$ with truncated training history (averaged across countries and horizons).}
\label{tab:short_history}
\small
\begin{tabular}{l RRR}
\toprule
Model & 20\,yr & 30\,yr & 50\,yr \\
\midrule
Frozen rates     & \textbf{0.296} & -- & \textbf{0.296} \\
LC-ResNet        & 0.317 & -- & 0.308 \\
Lee-Carter       & 0.322 & -- & 0.319 \\
LSTM $\kappa_t$  & 0.321 & -- & 0.317 \\
Random walk      & 0.328 & -- & 0.307 \\
Hyndman-Ullah    & 0.485 & -- & 0.480 \\
\bottomrule
\end{tabular}
\end{table}

\paragraph{Age groups.} LC-ResNet ranks first for young ages (0--19) with RMSE~0.310, suggesting that the neural residual correction captures non-linearities in infant and child mortality that the linear $b_x \kappa_t$ structure misses. For elderly ages (65--100), the random walk and Hyndman-Ullah perform best, while CBD---despite being designed for this segment---underperforms relative to models that benefit from information in younger ages.

% ══════════════════════════════════════════════════════════════
\section{Decision Framework}
\label{sec:decision}

Based on the benchmark results, we propose the following model selection guidelines:

\begin{table}[htbp]
\centering
\caption{Practical decision framework for model selection.}
\label{tab:decision}
\small
\begin{tabular}{p{4cm} l p{5cm}}
\toprule
Context & Recommended & Rationale \\
\midrule
Short history ($< 25$\,yr) & RWD & Neural networks lack training data; RWD is robust \\
Elderly ages (65+), regulatory & CBD & Designed for 60+; standard among regulators \\
Short horizon ($\le 5$\,yr) & LC-ResNet & Neural correction most active; best at $h=1$ \\
Medium horizon (5--10\,yr) & Poisson LC & Principled likelihood; stable \\
Long horizon ($> 10$\,yr) & Lee-Miller & Proven stability; $e_0$-adjusted \\
Interpretability required & LC-ResNet & LC skeleton is accessible; correction is bounded \\
\bottomrule
\end{tabular}
\end{table}

The key insight is that \textbf{no single model dominates across all scenarios}. The random walk is surprisingly competitive, and neural $\kappa_t$ replacements are only justified when the training history is long ($\ge 50$~years) and the forecast horizon is short ($\le 5$~years).

% ══════════════════════════════════════════════════════════════
\section{Actuarial Case Study}
\label{sec:casestudy}

We translate the benchmark results into financial terms using a stylized French annuity portfolio.

\subsection{Setup}

\begin{itemize}
    \item 1{,}000 pensioners aged 65
    \item Annual payment: 12{,}000~EUR per person
    \item Discount rate: $i = 2\%$ (risk-free)
    \item Mortality projection: each model's 10-year-ahead forecast
\end{itemize}

\subsection{Annuity pricing}

\begin{table}[htbp]
\centering
\caption{Annuity pricing results by model (French mortality data).}
\label{tab:pricing}
\small
\begin{tabular}{l r r r}
\toprule
Model & $\ddot{a}_{65}$ & Provision/person (EUR) & Total (M\,EUR) \\
\midrule
Lee-Carter  & 18.338 & 220{,}050 & 220.05 \\
Poisson LC  & 18.323 & 219{,}873 & 219.87 \\
LC-ResNet   & 18.268 & 219{,}217 & 219.22 \\
Random walk & 18.236 & 218{,}826 & 218.83 \\
CBD         & 17.816 & 213{,}793 & 213.79 \\
\midrule
\textbf{Model risk spread} & & & \textbf{6.26} \\
\bottomrule
\end{tabular}
\end{table}

The inter-model spread on total provisions reaches \textbf{6.26~million~EUR}---a material amount driven primarily by the CBD model, which projects higher mortality (lower $\ddot{a}_{65}$) for the elderly segment. Excluding CBD, the spread among full-grid models is 1.22~million~EUR.

\subsection{Solvency~II longevity shock}

Table~\ref{tab:scr} shows the impact of a standard Solvency~II longevity shock ($-20\%$ on mortality rates, i.e., policyholders live longer):

\begin{table}[htbp]
\centering
\caption{Solvency~II longevity shock analysis.}
\label{tab:scr}
\small
\begin{tabular}{l r r r}
\toprule
Model & Base (M\,EUR) & Shocked (M\,EUR) & SCR\textsubscript{long.} (M\,EUR) \\
\midrule
Lee-Carter  & 220.05 & 231.57 & 11.52 \\
Poisson LC  & 219.87 & 231.39 & 11.52 \\
LC-ResNet   & 219.22 & 230.80 & 11.58 \\
Random walk & 218.83 & 230.54 & 11.71 \\
CBD         & 213.79 & 226.12 & 12.33 \\
\midrule
\textbf{SCR range} & & & \textbf{0.81} \\
\bottomrule
\end{tabular}
\end{table}

The SCR\textsubscript{longevity} ranges from 11.52~M\,EUR (Lee-Carter) to 12.33~M\,EUR (CBD), a spread of 0.81~M\,EUR---approximately 7\% of the SCR itself. This demonstrates that model choice affects not only best-estimate provisions but also regulatory capital.

The SCR as a percentage of the base provision is remarkably stable across models (5.23--5.77\%), suggesting that the shock multiplier is robust to model specification, even if the base levels differ.

% ══════════════════════════════════════════════════════════════
\section{Discussion}
\label{sec:discussion}

\paragraph{The random walk puzzle.} The strongest finding of this benchmark is the competitiveness of the age-specific random walk with drift. It outperforms all classical LC variants and most neural models at horizons 5--20. This aligns with \citet{booth2006demographic} and challenges the assumption that structural models necessarily outperform na\"ive baselines.

\paragraph{Neural networks: promise and peril.} Neural $\kappa_t$ replacements (LSTM, GRU) show genuine improvement at short horizons ($h \le 5$), where the recurrent architecture captures recent trend changes that the linear drift misses. However, they degrade at long horizons, and this degradation is amplified on actuarial metrics because $e_0$ and $\ddot{a}_{65}$ are non-linear functions of the entire mortality surface. The Transformer and CNN architectures perform worst, likely due to limited training data relative to model complexity.

\paragraph{LC-ResNet: bounded flexibility.} The hybrid model achieves a favorable trade-off: it captures short-term residual structure (ranking first for young ages) while reverting to stable LC extrapolation at long horizons via the shrinkage mechanism. The interpretable LC skeleton ($a_x$, $b_x$, $\kappa_t$) remains accessible for regulatory reporting.

\paragraph{Limitations.}
\begin{itemize}
    \item All models are single-population; multi-population extensions (Li-Lee) are not considered.
    \item Cohort effects are not explicitly modeled (Renshaw-Haberman), though the LC-ResNet residual network may implicitly capture some.
    \item The neural models use a single seed (42) rather than ensemble averaging, which may underestimate their potential.
    \item HMD data ends in 2023; post-COVID stabilization effects remain uncertain.
\end{itemize}

% ══════════════════════════════════════════════════════════════
\section{Conclusion}
\label{sec:conclusion}

This benchmark provides an empirical answer to the question ``when should actuaries trust neural networks for mortality forecasting?'': at short horizons with long historical data, when the added complexity is justified by practical gains. For most actuarial applications---long-horizon projections for provisioning and capital---the classical Lee-Carter family and even the simple random walk remain competitive or superior.

The LC-ResNet hybrid offers a practical middle ground: interpretable by construction, competitive at short horizons, and stable at long horizons. The decision framework proposed in Section~\ref{sec:decision} provides actionable guidance for practitioners.

The full codebase---15~models implemented from scratch, evaluation harness, and Streamlit dashboard---is available at \url{https://github.com/aminemanai2003/neural-mortality-benchmark}.

% ══════════════════════════════════════════════════════════════
\bibliographystyle{apalike}
\bibliography{references}

\end{document}
